\section{Literature Survey}
We studied fifteen research papers in four areas that the project draws on: automated timetabling,
table recognition and OCR, natural-language interfaces to databases (NLIDB), and temporal reasoning
and security. Each is summarised below with what it contributes and what it leaves open for our
problem. Section~\ref{sec:gap} states the gap that our system fills.

\subsection{Literature Survey 1: A Survey of Automated Timetabling}
Schaerf~\cite{schaerf1999} classifies school, course and examination timetabling problems and the
search techniques used to solve them (graph colouring, integer programming, local search).
\textbf{Relevance:} it defines the hard clash constraints (no teacher, room or student group in two
places at once) that we \emph{check}. \textbf{Gap:} it assumes the timetable is generated by software;
reading an already published timetable is not addressed.

\subsection{Literature Survey 2: Recent Research Directions in Automated Timetabling}
Burke and Petrovic~\cite{burke2002} review constraint-based, meta-heuristic and hybrid approaches for
university timetabling and point out the distance between research prototypes and real institutional
data. \textbf{Relevance:} confirms that real data is messy and must drive the design. \textbf{Gap:} no
support for the consumers of a timetable after it is published.

\subsection{Literature Survey 3: Meta-heuristics for University Timetabling}
Lewis~\cite{lewis2008} surveys meta-heuristic techniques and separates hard constraints from soft ones
(preferences, compactness). \textbf{Relevance:} the hard constraints become our conflict rules.
\textbf{Gap:} construction only; no search, role views or clash checking of a published grid.

\subsection{Literature Survey 4: A Survey of Table Recognition}
Zanibbi, Blostein and Cordy~\cite{zanibbi2004} describe table recognition as detection, structure
recognition (rows, columns, spanning cells) and interpretation, and note that interpretation is domain
specific. \textbf{Relevance:} motivated our timetable-specific grid parser. \textbf{Gap:} no
interpretation of lab batches, legends or time rows.

\subsection{Literature Survey 5: The T-Recs Table Recognition System}
Kieninger and Dengel~\cite{kieninger1998} recognise tables from word bounding boxes without relying on
ruling lines. \textbf{Relevance:} basis of our geometric fallback that rebuilds columns from whitespace
gutters when a page has no ruled lines. \textbf{Gap:} generic tables only.

\subsection{Literature Survey 6: ICDAR 2013 Table Competition}
G\"obel et~al.~\cite{gobel2013} define standard metrics for table location and structure recognition
and show that cell-structure recovery remains error-prone. \textbf{Relevance:} led us to give every
extracted entry a verification status instead of trusting the extractor. \textbf{Gap:} evaluation of
extraction, not of downstream use.

\subsection{Literature Survey 7: An Overview of the Tesseract OCR Engine}
Smith~\cite{smith2007} describes Tesseract's line finding, adaptive classifier and word recognition.
\textbf{Relevance:} we use Tesseract for scanned pages and photographs. \textbf{Gap:} OCR output has
no table semantics, so OCR entries are capped at \textsc{Unverified} in our system.

\subsection{Literature Survey 8: Natural Language Interfaces to Databases}
Androutsopoulos, Ritchie and Thanisch~\cite{androutsopoulos1995} review pattern-matching, syntax-based
and semantic-grammar NLIDB systems and observe that a narrow domain allows simple, reliable
rule-based interpretation. \textbf{Relevance:} supports our deterministic planner. \textbf{Gap:}
classical systems map a question to one fixed pattern; we compose independent pieces into a plan.

\subsection{Literature Survey 9: Seq2SQL}
Zhong, Xiong and Socher~\cite{zhong2017} generate SQL from text with a neural sequence model trained
by reinforcement learning. \textbf{Relevance:} shows the flexibility of learned text-to-SQL.
\textbf{Gap:} needs training data and is not exact; a wrong query sends a student to the wrong room.

\subsection{Literature Survey 10: The Spider Text-to-SQL Benchmark}
Yu et~al.~\cite{yu2018} introduce a large cross-domain text-to-SQL benchmark; exact-match accuracy on
complex queries is far below 100\%. \textbf{Relevance:} quantifies the accuracy risk of generated SQL.
\textbf{Gap:} no guarantee of correctness or safety for a single deployed domain.

\subsection{Literature Survey 11: Indirect Prompt Injection}
Greshake et~al.~\cite{greshake2023} show that LLM-integrated applications can be taken over by
instructions hidden in the data they read. \textbf{Relevance:} letting an LLM drive data access is a
security risk. \textbf{Gap:} mitigation remains an open problem.

\subsection{Literature Survey 12: From Prompt Injection to SQL Injection}
Pedro et~al.~\cite{pedro2023} demonstrate prompts that make LLM web applications run malicious SQL.
\textbf{Relevance:} the main reason our engine never builds SQL from user text and uses only
pre-written parameterised templates.

\subsection{Literature Survey 13: Maintaining Knowledge about Temporal Intervals}
Allen~\cite{allen1983} defines thirteen relations between time intervals. \textbf{Relevance:} our clash
rule uses the overlap family $a.s < b.e \wedge b.s < a.e$ with end-exclusive intervals.

\subsection{Literature Survey 14: Role-Based Access Control Models}
Sandhu et~al.~\cite{sandhu1996} formalise RBAC with roles, permissions and sessions.
\textbf{Relevance:} our five roles and their data scopes, enforced on the server for every request.

\subsection{Literature Survey 15: Argon2 Memory-Hard Hashing}
Biryukov, Dinu and Khovratovich~\cite{biryukov2016} present Argon2, the winner of the Password Hashing
Competition, which resists GPU cracking. \textbf{Relevance:} used for all stored passwords.

\subsection{Comparison with Existing Systems}
Besides the papers, we compared the tools a college could use today (Table~\ref{tab:compare}):
timetable generators (FET, UniTime, aSc~\cite{fet,unitime,asc}), PDF table extractors (Tabula,
Camelot~\cite{tabula,camelot}), calendar apps, ERP modules and LLM text-to-SQL.
\begin{table}[H]
\centering
\caption{Comparison of existing approaches with the proposed system}
\label{tab:compare}
\footnotesize
\setlength{\tabcolsep}{3pt}
\begin{tabular}{@{}p{3.6cm}*{7}{>{\centering\arraybackslash}p{1.6cm}}@{}}
\toprule
\textbf{Criterion} & \textbf{Manual PDF} & \textbf{FET / UniTime / aSc} & \textbf{Tabula / Camelot} & \textbf{Calendar apps} & \textbf{ERP module} & \textbf{LLM text-to-SQL} & \textbf{Proposed}\\
\midrule
Reads published PDF/photo & -- & No & Partly & No & No & No & \textbf{Yes}\\
Timetable semantics (labs, legends) & Human & Yes & No & No & Yes & Depends & \textbf{Yes}\\
Role-based views & No & Limited & No & No & Yes & No & \textbf{Yes (5)}\\
Plain-English questions & No & No & No & Limited & No & Yes & \textbf{Yes (composed plans)}\\
Reproducible, injection-safe answers & -- & -- & -- & -- & Yes & No & \textbf{Yes}\\
Clash detection & No & At creation & No & No & Sometimes & No & \textbf{Yes, after publication}\\
Version archive & Ad hoc & Files & No & No & Sometimes & No & \textbf{Yes, audited}\\
Personal timetable kept private & -- & No & No & Yes & No & No & \textbf{Yes}\\
Cost & Free & Free/Paid & Free & Free & Paid & Paid API & \textbf{Free (open source)}\\
Mobile use & PDF viewer & Limited & No & Yes & Varies & Varies & \textbf{Android, iOS, web}\\
\bottomrule
\end{tabular}
\end{table}


\subsection{Research Gap}\label{sec:gap}
\begin{enumerate}
\item Timetabling research and generators \emph{construct} timetables; none reads the timetable as it
      is published or helps the people who follow it~\cite{schaerf1999,burke2002,lewis2008}.
\item Table recognition recovers cells but not timetable meaning: stacked lab batches, page-wise
      legends and mislabelled time rows~\cite{zanibbi2004,gobel2013}.
\item Text-to-SQL and LLM interfaces are flexible but not exact and are open to
      injection~\cite{yu2018,greshake2023,pedro2023}.
\item No tool used by our college checks the \emph{published} timetable for clashes.
\end{enumerate}
The gap is therefore not a better construction algorithm, but the missing link between the published
document and its daily users: trustworthy, role-scoped, searchable data with provenance, clash checks
and version history, and safe answers to plain-English questions.
\newpage
